\section{总结与延伸}

\subsection{核心结论}

\begin{enumerate}
\item Agent 时代不是聊天能力增强，而是模型进入环境、工具和长程任务。
\item OpenClaw 的价值在于触发框架共识，让团队看到模型与 scaffold 互相放大的可能。
\item 后训练尤其是 Agent RL scaling 成为主线，需要任务环境、rollout、reward、eval 和成本控制。
\item MiMo-V2 的意义在于多模型编排，而不是单个模型发布。
\item 算力分配从预训练独大转向 research、pre-train、post-train 更均衡。
\item 组织平权、多样性和文化价值观会影响 AI Lab 拥抱新范式的速度。
\item 国内团队的追赶机会在 Agent post-train、成本敏感和应用反馈，但挑战在 RL infra 和组织敏捷性。
\end{enumerate}

\subsection{开放问题}

\begin{itemize}
\item Agent post-train 的主流 reward 和环境形态会是什么？
\item 1T 基座模型是否长期是入场券，还是会被小模型 + 框架逐步侵蚀？
\item 多模型编排会走向统一生态，还是开放市场中的模型路由？
\item 组织平权如何避免变成混乱协作，而真正提升研究速度？
\item 当模型没有“生存危机”时，人类应该如何设计它的价值函数和社会位置？
\end{itemize}

\subsection{拓展阅读}

\begin{itemize}
\item EP139 苏煜 Agent 技术史：理解 OpenClaw Moment 与 Language Agent 的历史背景。
\item Claude Code / OpenClaw / Computer Use Agent 相关材料：理解 Agent 框架如何改变产品形态。
\item RLHF、RLAIF、verifier、rollout、agent benchmark 相关论文：理解后训练基础设施。
\item MoE、attention alternatives、multi-modal model orchestration：理解 MiMo-V2 背后的模型系统问题。
\end{itemize}

\begin{importantbox}{最后的判断}
Agent 时代的胜负不会由某一个单点模型决定。它更像一场系统竞赛：谁能更快地把模型、框架、任务环境、后训练、算力、产品和组织拧成闭环，谁就更可能在下一阶段领先。
\end{importantbox}

\end{document}
